\chapter{Implementation}
\label{ch:implementation}
\vspace{2em}

% TODO: Introduce the chapter -- describe how the system design from
% Chapter~\ref{ch:design} was realised in practice.

\section{Technology Stack}
\label{impl:sec:stack}

% TODO: Languages, frameworks, models, and infrastructure used.

\section{Knowledge Base and Retrieval}
\label{impl:sec:retrieval}

% TODO: How course materials and administrative data were ingested,
% indexed, and retrieved.

\section{Question Answering and Personalisation}
\label{impl:sec:qa}

% TODO: The LLM integration, prompting/RAG strategy, and personalisation
% logic.

\section{Context Compression for Token Reduction}
\label{impl:sec:compression}

A practical constraint of the deployed system is that response generation runs on a
small, locally hosted language model (Llama~3.2, 3B parameters, served through
Ollama on an 8\,GB consumer laptop). Under standard retrieval-augmented generation
(RAG)~\cite{neurips20:Lewis}, the top-$N$ retrieved chunks of course material are
prepended to the prompt in full, even though each chunk is typically only
partially relevant to the student's question. This inflates prompt length --- the
dominant cost for a compute-bound local model --- and risks distracting the model
with irrelevant context, a failure mode documented for long-context
inputs~\cite{tacl24:Liu}.

To address this, the pipeline implements a lightweight context-compression stage
inspired by \textsc{Recomp}~\cite{iclr24:Xu}, which proposes compressing retrieved
passages before generation and, crucially, \emph{selective augmentation}:
returning an empty context when nothing relevant survives compression, rather than
forcing the model to attend to noise. The implementation here is the untrained
extractive variant: instead of the paper's fine-tuned contrastive compressor, it
selects existing sentences by embedding similarity, making it faithful by
construction (the compressor cannot introduce content that was not retrieved).

The stage is inserted between the existing retriever and the generation call, and
proceeds as follows:

\begin{enumerate}[nosep]
  \item retrieved chunks are segmented into candidate sentences, with markdown
        structure and code fences stripped;
  \item each candidate sentence is embedded with the system's existing retrieval
        embedder (BGE-M3~\cite{arxiv24:Chen}) in a single batched call;
  \item each sentence is scored by cosine similarity against the query embedding,
        which is reused from the retrieval step at no additional cost;
  \item sentences scoring above a threshold $\tau$ are retained, highest first,
        capped at $k$ sentences; if none qualifies, the context is left empty
        (selective augmentation) and the model is instructed to state that the
        topic is not covered by the course materials.
\end{enumerate}

Reusing the retrieval embedder is a deliberate engineering choice: the compressor
adds no new model, dependency, or runtime to the stack (approximately 130 lines of
JavaScript), which matters on the memory-constrained host described in
\autoref{impl:sec:stack}. The stage is deployed behind a runtime feature flag
(\texttt{RAG\_COMPRESS}, with $\tau$ and $k$ configurable), which both protects
the production path and enables the controlled A/B comparison reported in
\autoref{eval:sec:results}.

\section{Analytics and Feedback Loop}
\label{impl:sec:analytics}

% TODO: How student-question analytics and performance insights are
% surfaced to faculty.

\section{Summary}
\label{impl:sec:summary}

% TODO: Recap what was built, leading into the evaluation.
